% Compile from the FSI repository root:
%   docs/resources/benchmarks/abramson1965_report/build_chapter_close_report.sh
\documentclass[11pt,a4paper]{article}

\usepackage[utf8]{inputenc}
\usepackage[T1]{fontenc}
\usepackage[italian]{babel}
\usepackage{geometry}
\usepackage{graphicx}
\usepackage{subcaption}
\usepackage{booktabs}
\usepackage{siunitx}
\usepackage{amsmath}
\usepackage{amssymb}
\usepackage{float}
\usepackage{xcolor}
\usepackage{hyperref}
\usepackage{enumitem}

\makeatletter
\newcommand{\fsiroot}{}
\newcommand{\fsipath}[1]{\fsiroot#1}
\IfFileExists{output_data/ABRAMSON1965/vlambda_data.csv}{%
  \def\fsiroot{}%
}{%
  \IfFileExists{../../../../output_data/ABRAMSON1965/vlambda_data.csv}{%
    \def\fsiroot{../../../../}%
  }{%
    \PackageWarning{abramson1965}{Cannot locate FSI repo root; figures/CSV may be missing.}%
  }%
}
\makeatother

\geometry{margin=22mm}
\sisetup{detect-all=true, per-mode=symbol}
\DeclareSIUnit{\knunit}{knots}
\DeclareSIUnit{\cps}{cps}
\DeclareSIUnit{\ftunit}{ft}

\title{ABRAMSON1965 --- chiusura capitolo\\[0.4em]
\large Riproduzione idrodinamica, accoppiamento aero--struttura e campagna di validazione}
\author{Report tecnico di lavoro --- FSI / lorebasket}
\date{23 giugno 2026}

\begin{document}
\maketitle

\begin{abstract}
Questo documento chiude il capitolo di lavoro sul benchmark idro-aeroelastico
ABRAMSON1965 (Abramson \& Ransleben, 1965).
Si riassume la \textbf{pipeline numerica} costruita (beam FE, DLM/PanelAero, Roger RFA,
Capytaine BEM, strip theory, tracking modale), la \textbf{campagna di test} salvata in
\texttt{output\_data/ABRAMSON1965/}, il \textbf{punto di arrivo} sulla riproduzione
dell'idrodinamica e le \textbf{limitazioni} residue rispetto al flutter sperimentale
(\SI{48.1}{\knunit}, $f\approx\SI{17.5}{\cps}$).
La linea di riferimento operativa e' \texttt{RFA\_PK\_method} con massa aggiunta
Capytaine in $M_\mathrm{eff}$; il gap dominante sul flutter resta lo
smorzamento circolatorio DLM ($\mathrm{Im}(Q)$), non la sola massa aggiunta.
\end{abstract}

\tableofcontents
\newpage

%==============================================================================
\section{Riferimento sperimentale e obiettivo}
%==============================================================================

Il caso Abramson e' un idrofoil NACA 16-012 completamente immerso, cantilever,
con mass ratio $\mu\approx 1$ e frequenze a secco $\approx\SI{10.5}{\hertz}$ (flessione)
e $\approx\SI{20.6}{\hertz}$ (torsione).
La rottura per flutter avviene a \textbf{\SI{48.1}{\knunit}} con frequenza accoppiata
$\approx\SI{17.5}{\cps}$ e $(V/b\omega)_F\approx 1{,}48$.

\paragraph{Obiettivo del capitolo.}
Costruire una catena numerica riproducibile che:
\begin{enumerate}[nosep]
  \item riproduca modi secchi e parametri strutturali del paper;
  \item integri idrodinamica potenziale (massa aggiunta e, opzionalmente, smorzamento radiativo)
        via Capytaine o strip theory;
  \item accoppi aerodinamica non stazionaria DLM in spazio modale $Q_\mathrm{modal}$;
  \item produca diagrammi $V$--$g$, $V$--$f$, $V$--$\delta$ confrontabili con le Figure~5--6;
  \item quantifichi sensibilita' a massa aggiunta, allineamento $x_{EA}$/aerogrid e
        al fattore \texttt{aero\_im\_Q\_scale}.
\end{enumerate}

%==============================================================================
\section{Architettura della pipeline}
%==============================================================================

\subsection{Flusso di calcolo}

\begin{enumerate}[nosep]
  \item \textbf{Struttura a secco:} beam 100 elementi $\to$ modi $\Phi$, $\hat{M},\hat{C},\hat{K}$,
        Rayleigh $\zeta=0{,}03$ sui primi due modi.
  \item \textbf{Idrodinamica:} Capytaine BEM su mesh VTU (4 modi) \emph{oppure} strip theory
        proiettata su $\Phi$; $A_\infty$, $B_\infty$ in $M_\mathrm{eff}$, $C_\mathrm{eff}$.
  \item \textbf{Aerodinamica:} $Q_{jj}(k)$ DLM precomputata; $Q_\mathrm{modal}(k)$ via
        $Z$, $Z_{qs}$, $Z_\mathrm{force}$, $\Phi$.
  \item \textbf{Flutter:} Roger RFA ($B_0,B_1,B_2,\,B_\mathrm{lag}$) + autovalori del sistema
        aumentato; tracking MAC su due rami (bend/pitch).
\end{enumerate}

\subsection{Convenzione del solver (\texttt{RFA\_PK})}

\[
  M_\mathrm{eff} = \hat{M} + A_\infty,
  \qquad
  K_\mathrm{eff} = \hat{K} - \tfrac{1}{2}\rho V^2 B_0,
  \qquad
  C_\mathrm{eff} = \hat{C} - s_Q \tfrac{1}{2}\rho V B_1 + B_\infty + C_\mathrm{emp},
\]
con $s_Q=\texttt{aero\_im\_Q\_scale}$.
La parte immaginaria di $Q$ alimenta $B_1$ e quindi lo smorzamento aerodinamico;
$\mathrm{Re}(Q)$ entra in $B_0$ (rigidita' aerodinamica).
Nel blocco bend--pitch accoppiato, \textbf{non} basta il segno delle diagonali
$\mathrm{Re}(Q_{\alpha\alpha})$: le frequenze dei rami sono autovalori di un sistema
$2\times2$ e dipendono anche da $\mathrm{Re}(Q_{h\alpha})$, $\mathrm{Re}(Q_{\alpha h})$.

\subsection{Configurazione nominale}

File: \texttt{examples/ABRAMSON1965/config\_ABRAMSON1965.py}.
Parametri chiave: $x_{EA}/c=0{,}25$, $x_{CM}/c=0{,}512$, profondita' \SI{0.0508}{\meter},
\texttt{num\_modes\_egv}=4, \texttt{num\_modes\_flutter\_egv}=2,
\texttt{RFA\_PK\_method=True}, \texttt{rfa\_apply\_b2\_mass=False}.

%==============================================================================
\section{Campagna di test (\texttt{output\_data/ABRAMSON1965})}
%==============================================================================

Tutti i risultati sono organizzati sotto \texttt{output\_data/ABRAMSON1965/}.
Ogni cartella contiene tipicamente \texttt{vg\_data.csv}, \texttt{vf\_data.csv},
\texttt{vlambda\_data.csv}, \texttt{flutter\_summary.csv}, figure PNG e \texttt{run.log}.

\subsection{Panoramica delle directory}

\begin{table}[H]
  \centering
  \caption{Macro-aree della campagna numerica.}
  \small
  \begin{tabular}{p{3.6cm}p{9.8cm}}
    \toprule
    Directory & Contenuto \\
    \midrule
    \texttt{00\_baseline\_sensitivity} &
    Confronto strip / Capytaine BEM / Capytaine striplike (stessa DLM, solo $M_\mathrm{hydro}$ cambia). \\
    \texttt{01\_uncoupled\_rfa\_capytaine} &
    Riferimento nominale $x_{EA}=x_{CM}$; aerogrid fissa; Capytaine BEM. \\
    \texttt{02\_uncoupled\_xea\_xcm\_sweep} &
    Sweep allineato $x_{EA}=x_{CM}\in\{0.20,0.22,0.38,0.40\}$; aerogrid decoupled a $0{,}25c$. \\
    \texttt{alignment\_sweep} &
    Sweep $x_{EA}$ con Capytaine per caso; varianti strip e DLM puro. \\
    \texttt{input\_sensitivity} &
    Perturbazioni singole ($\mu$, $EI$, $GJ$, $x_{EA}$, $\zeta$, scala strip$\to$Capytaine). \\
    \texttt{measured\_q\_tests} &
    Ibridi $\mathrm{Re}(Q)_\mathrm{DLM}+\mathrm{Im}(Q)_\mathrm{misurato}$, scale $\mathrm{Im}(Q)$. \\
    \texttt{general\_results} & Run storiche PK/RFA+strip, scale \texttt{aero\_im\_Q\_scale}. \\
    \bottomrule
  \end{tabular}
\end{table}

\subsection{Baseline idroelastico (cartella \texttt{00\_baseline\_sensitivity})}

Tre run con \textbf{stessa aerodinamica DLM} e tracking su due rami; cambia solo la
proiezione della massa aggiunta in $M_\mathrm{hat}$:

\begin{table}[H]
  \centering
  \caption{Flutter baseline RFA+PK (Capytaine in $M_\mathrm{eff}$, $\zeta=0{,}03$, $s_Q=1$).}
  \label{tab:baseline-flutter}
  \begin{tabular}{lcccl}
    \toprule
    Run & $V_F$ [kn] & Ramo & $f_F$ [Hz] & Modo instabile \\
    \midrule
  Sperimentale & 48.1 & --- & 17.5 & --- \\
  \texttt{\_baseline\_rfa\_strip} & 32.4 & 1 & 11.9 & bend-like \\
  \texttt{\_baseline\_rfa\_capytaine\_striplike} & 30.8 & 1 & 12.0 & bend-like \\
  \texttt{\_baseline\_rfa\_capytaine} (BEM) & 37.7 & 2 & 11.8 & torsion-like \\
    \bottomrule
  \end{tabular}
\end{table}

\paragraph{Interpretazione.}
La sola modifica della \textbf{distribuzione modale} di $A_\infty$ sposta quale ramo
va in flutter per primo: con strip/striplike domina il bend (ramo~1); con Capytaine BEM
($A_{22}\approx\SI{2.0}{\kilogram}$ sul pitch, $\approx 2{,}3\times$ strip) il pitch
e' fortemente inerziale e fluttera il ramo~2.
Non e' un cambio di $Q_{jj}$: e' \textbf{riorganizzazione dello spettro accoppiato}
bend--torsione e del tracking MAC.

\begin{figure}[H]
  \centering
  \begin{subfigure}{0.32\textwidth}
    \centering
    \includegraphics[width=\linewidth]{\fsipath{output_data/ABRAMSON1965/00_baseline_sensitivity/_baseline_rfa_strip/vg_vf_combined.png}}
    \caption{Strip}
  \end{subfigure}\hfill
  \begin{subfigure}{0.32\textwidth}
    \centering
    \includegraphics[width=\linewidth]{\fsipath{output_data/ABRAMSON1965/00_baseline_sensitivity/_baseline_rfa_capytaine_striplike/vg_vf_combined.png}}
    \caption{Striplike}
  \end{subfigure}\hfill
  \begin{subfigure}{0.32\textwidth}
    \centering
    \includegraphics[width=\linewidth]{\fsipath{output_data/ABRAMSON1965/00_baseline_sensitivity/_baseline_rfa_capytaine/vg_vf_combined.png}}
    \caption{Capytaine BEM}
  \end{subfigure}
  \caption{Baseline sensitivity: stesso solver, tre proiezioni di massa aggiunta.}
  \label{fig:baseline-trio}
\end{figure}

\subsection{Massa aggiunta modale: strip, striplike, BEM}

\begin{table}[H]
  \centering
  \caption{Diagonali $A_\infty$ modali [kg] sui primi quattro modi secchi
  (\texttt{added\_mass\_projection\_comparison.json}, caso striplike).}
  \label{tab:added-mass}
  \begin{tabular}{lcccc}
    \toprule
    Fonte & Modo 0 & Modo 1 & Modo 2 & Modo 3 \\
    \midrule
    Strip theory & 0.516 & 0.878 & 0.140 & 0.186 \\
    Capytaine striplike & 0.414 & 0.706 & 0.113 & 0.149 \\
    Capytaine BEM & 0.198 & \textbf{1.994} & 0.374 & 0.363 \\
    \midrule
    BEM / strip (modo 1) & 0.38 & \textbf{2.27} & 2.67 & 1.95 \\
    \bottomrule
  \end{tabular}
\end{table}

Il BEM 3D \textbf{abbassa} la massa aggiunta sul bend (modo~0) ma la \textbf{moltiplica}
sul pitch (modo~1), con accoppiamenti off-diagonal non trascurabili
($A_{12}\approx\SI{-0.48}{\kilogram}$ nella documentazione Capytaine).
La proiezione striplike e' una via di mezzo fisicamente motivata per modi flesso-torsionali
quando il BEM puro sovrastima l'inerzia torsionale.

\begin{figure}[H]
  \centering
  \includegraphics[width=0.88\linewidth]{\fsipath{output_data/ABRAMSON1965/ABRAMSON1965_Madded_terms_vs_V_capytaine_vs_strip.png}}
  \caption{Confronto termini di massa aggiunta modale (strip vs Capytaine).}
\end{figure}

\subsection{Sweep $x_{EA}=x_{CM}$ decoupled (\texttt{02\_uncoupled\_xea\_xcm\_sweep})}

Quattro casi con aerogrid DLM fissa a $x_{EA}/c=0{,}25$ e barra strutturale spostata:
\begin{itemize}[nosep]
  \item \texttt{xea\_0p20}, \texttt{xea\_0p22} (verso prua);
  \item \texttt{xea\_0p38}, \texttt{xea\_0p40} (verso poppa).
\end{itemize}
Per ogni caso: modi secchi $\to$ Capytaine BEM ricalcolato $\to$ RFA\_PK flutter.
Con \texttt{num\_modes\_egv}=4 (pool di 4 radici, 2 rami tracciati), tutti e quattro i casi
presentano \textbf{0/120 NaN} su Branch~2 in \texttt{vf\_data} (tracking continuo su tutto
lo sweep fino a $V_\mathrm{max}=\SI{40}{\knunit}$). Nessun caso mostra attraversamento
$g=0$ nel range di velocita' esplorato.

\begin{figure}[H]
  \centering
  \includegraphics[width=0.95\linewidth]{\fsipath{output_data/ABRAMSON1965/02_uncoupled_xea_xcm_sweep/uncoupled_xea_xcm_sweep_panel.png}}
  \caption{Pannello sweep $x_{EA}=x_{CM}$ con \texttt{num\_modes\_egv}=4: rami bend/pitch
  continui su tutti i casi (prua e poppa).}
\end{figure}

\subsection{Altre campagne rilevanti}

\begin{itemize}
  \item \textbf{\texttt{input\_sensitivity/strip\_added\_mass\_scale}:}
        interpolazione strip$\to$Capytaine; $V_F$ scende da $\sim\SI{33}{\knunit}$ (strip puro)
        a $\sim\SI{24}{\knunit}$ (Capytaine pieno) --- coerente con maggiore inerzia sul pitch.
  \item \textbf{\texttt{alignment\_sweep/xea\_sweep\_capytaine}:}
        sweep fine verso poppa con Capytaine per caso; coalescenza modi $\Rightarrow$ criticita'
        di tracking (PAS, \texttt{rfa\_pas\_allow\_coalescence}).
  \item \textbf{\texttt{measured\_q\_tests}:}
        sostituzione parziale di $\mathrm{Im}(Q)$ con coefficienti misurati Abramson;
        scala \texttt{aero\_im\_Q\_scale} come knob fenomenologico.
  \item \textbf{Report precedente} (\texttt{abramson1965\_current\_report.tex}):
        confronto $\mathrm{Im}(Q_{bb})$ DLM vs $L_{hh}$ misurato ($\times 7$--$\times 12$ a basse $V$).
\end{itemize}

%==============================================================================
\section{Fisica dell'accoppiamento e del flutter numerico}
%==============================================================================

\subsection{Ruolo di $\mathrm{Re}(Q)$ e $\mathrm{Im}(Q)$}

\begin{itemize}[nosep]
  \item $\mathrm{Re}(Q_{\alpha\alpha})>0$: softening aerodinamico $\Rightarrow$ tendenza a
        \emph{abbassare} $\omega_\mathrm{wet}$ (non equivale a flutter).
  \item $\mathrm{Im}(Q)$: smorzamento circolatorio; attraversamento $\gamma=0$ $\Leftrightarrow$ flutter.
  \item Off-diagonal $\mathrm{Re}(Q_{h\alpha})$, $\mathrm{Im}(Q_{h\alpha})$: coalescenza
        bend--torsione; determinano quale combinazione modale diventa instabile.
\end{itemize}

\subsection{Paradosso ``non conservativo'' vs flutter anticipato}

La letteratura inviscida spesso \emph{sovrastima} $V_F$ (manca smorzamento viscoso stabilizzante).
Nel codice il flutter e' \textbf{anticipato} perch\'e:
\begin{enumerate}[nosep]
  \item $\mathrm{Im}(Q)_\mathrm{DLM}$ in acqua ($\rho\approx 997$) genera $C_\mathrm{aero}$
        molto piu' grande di $C_\mathrm{struct}$ ($\zeta=3\%$);
  \item $C_\mathrm{hydro}\approx 0$ nella banda di flutter ($B_\infty\to 0$);
  \item il paper usa coefficienti \textbf{misurati} (curva~3), non DLM puro.
\end{enumerate}
Con \texttt{aero\_im\_Q\_scale} $=0{,}5$ e $0{,}4$ si recupera rispettivamente
$V_F\approx\SI{18}{\knunit}$ e $\approx\SI{31}{\knunit}$ (RFA+strip storico),
avvicinando $f_F$ al valore sperimentale ma ancora lontano da \SI{48.1}{\knunit}.

\begin{figure}[H]
  \centering
  \includegraphics[width=0.92\linewidth]{\fsipath{output_data/ABRAMSON1965/00_baseline_sensitivity/_baseline_rfa_capytaine_striplike/stiffness_damping_contributions.png}}
  \caption{Contributi $M/K/C$: dominanza di $C_\mathrm{aero}$ (B1) sul modo bend a $V$ crescenti.}
\end{figure}

%==============================================================================
\section{Tracking modale e limiti numerici}
%==============================================================================

\subsection{Problema \texttt{mode\_lost} (pool a 2 radici)}

Con \texttt{num\_modes\_egv}=2 il pool di autovalori oscillatori contiene solo due radici.
Durante la coalescenza bend--pitch ($V\sim\SI{17}{\knunit}$), la radice pitch esce dal
top-2 e ne entra una spuria ($\omega\sim\SI{36}{\radian\per\second}$, $\gamma\ll 0$):
il MAC non trova corrispondenza $\Rightarrow$ Branch~2 \texttt{LOST} per $\sim$20 step
$\Rightarrow$ buchi nei CSV \texttt{vf\_data}/\texttt{vg\_data}.

\paragraph{Correzione adottata.}
\texttt{num\_modes\_egv}=4 (pool di ricerca a 4 radici), \texttt{num\_modes\_flutter\_egv}=2
(ancora due rami in output). Rerun completati su tutti e quattro i casi
(\texttt{xea\_0p20}--\texttt{xea\_0p40}): \textbf{0/120 buchi} su Branch~2.

\subsection{Altre limitazioni strutturali del modello}

\begin{enumerate}
  \item \textbf{Geometria DLM fissa:} nessun aggiornamento di $Q_{jj}$ con deformazione;
        valido solo per piccole oscillazioni.
  \item \textbf{Decoupling aerogrid:} negli sweep \texttt{02\_*} la griglia DLM resta a
        $x_{EA}/c=0{,}25$ mentre la trave si sposta --- utile per sensitivita' strutturale,
        non fisicamente coerente per confronto assoluto con Abramson.
  \item \textbf{Capytaine:} profondita' finita, superficie libera a $z=0$; mesh resolution
        warning a $\omega$ alte; $B(\omega)\approx 0$ in banda flutter.
  \item \textbf{Roger RFA:} fit su griglia $k$ con errori residui fino a $\sim 30\%$ su
        $\mathrm{Re}(Q)$; lag terms fino a $\sim 18\times \|B_0\|$.
  \item \textbf{Proxy $\mathrm{Im}(Q)$ vs misurato:} confronto su $Q_{bb}$ vs $L_{hh}$ scalato,
        non equivalenza strip completa del paper.
  \item \textbf{Flutter ancora anticipato:} anche la migliore baseline Capytaine
        ($\sim\SI{38}{\knunit}$) resta $\sim 20\%$ sotto il test.
\end{enumerate}

%==============================================================================
\section{Punto di arrivo: cosa e' stato riprodotto}
%==============================================================================

\begin{table}[H]
  \centering
  \caption{Stato di riproduzione rispetto ad Abramson (1965).}
  \label{tab:reproduction-status}
  \begin{tabular}{p{4.2cm}cc}
    \toprule
    Aspetto & Stato & Nota \\
    \midrule
    Geometria e masse Tab.~1 & \textcolor{teal}{OK} & $b$, $L$, $\mu$, $EI$, $GJ$, $r_\alpha^2$ \\
    Frequenze a secco & \textcolor{teal}{OK} & $\approx\SI{10.5}{Hz}$, $\approx\SI{20.6}{Hz}$ \\
    Pipeline automatizzata & \textcolor{teal}{OK} & dry $\to$ Capytaine $\to$ RFA\_PK $\to$ post \\
    Massa aggiunta 3 vie & \textcolor{teal}{OK} & strip / striplike / BEM confrontabili \\
    Smorzamento radiativo $B(\omega)$ & \textcolor{teal}{OK} & calcolato; $\approx 0$ a flutter \\
    Trend $f(V)$ coalescenza & \textcolor{teal}{OK} & qualitativo come Fig.~6 \\
    $V_F$ assoluto & \textcolor{red}{Gap} & $\sim 31$--$38$ kn vs 48.1 kn \\
    $\mathrm{Im}(Q)$ vs misurato & \textcolor{orange}{Parziale} & proxy $L_{hh}$; gap $\times 7$--$12$ \\
    Coefficienti misurati curva~3 & \textcolor{orange}{Parziale} & test ibridi, non baseline \\
    Viscosita' / superficie libera & \textcolor{red}{Assente} & DLM inviscido \\
    \bottomrule
  \end{tabular}
\end{table}

\paragraph{Configurazione raccomandata per studi successivi.}
\begin{itemize}[nosep]
  \item \texttt{RFA\_PK\_method=True}, Capytaine BEM o striplike a seconda del trade-off
        inerzia torsionale;
  \item \texttt{num\_modes\_egv=4}, \texttt{num\_modes\_flutter\_egv=2};
  \item \texttt{rfa\_branch\_guard\_enable=True}, policy \texttt{stable};
  \item per allineamento $V_F$: calibrare \texttt{aero\_im\_Q\_scale} e/o ibride
        $\mathrm{Im}(Q)_\mathrm{misurato}$; non aspettarsi match con solo $A_\infty$.
\end{itemize}

%==============================================================================
\section{Conclusioni}
%==============================================================================

Il capitolo ABRAMSON1965 ha portato a una \textbf{piattaforma idro-aeroelastica completa}:
modi secchi verificati, idrodinamica potenziale da Capytaine con proiezione strip/striplike,
accoppiamento DLM modale, solver Roger RFA con tracking MAC, e una campagna estesa di
sensibilita' archiviata sotto \texttt{output\_data/ABRAMSON1965/}.

Sul piano \textbf{idrodinamico}, la massa aggiunta modale e' stata riprodotta e confrontata
su tre livelli; l'effetto sulla dinamica accoppiata e' \textbf{forte e asimmetrico} tra
bend e pitch, spiegando l'inversione del ramo flutterante tra strip e Capytaine BEM.

Sul piano \textbf{aerodinamico}, il gap residuo verso \SI{48.1}{\knunit} e' dominato da
$\mathrm{Im}(Q)_\mathrm{DLM}$ troppo grande rispetto ai coefficienti oscillatori misurati,
non da errori nella sola $A_\infty$.
La correzione fenomenologica \texttt{aero\_im\_Q\_scale} e le matrici misurate sono le
levette calibrate disponibili nel codice.

Sul piano \textbf{numerico}, il tracking con pool a due radici era insufficiente in
coalescenza; \texttt{num\_modes\_egv=4} risolve i buchi sui diagrammi $V$--$f$.

\medskip
\noindent\textbf{Chiusura.}
Il benchmark e' \emph{riproducibile} e \emph{diagnosticabile}; non e' ancora \emph{predittivo}
a $V_F$ sperimentale senza calibrazione di $\mathrm{Im}(Q)$.
Il report tecnico di dettaglio su input, $\delta(V)$ e confronto Fig.~5--6 resta
\texttt{abramson1965\_current\_report.tex}; la sensibilita' parametrica in
\texttt{sensitivity\_rerun\_report.tex}.

\appendix
\section{Script e percorsi utili}

\begin{itemize}[nosep]
  \item Config: \texttt{examples/ABRAMSON1965/config\_ABRAMSON1965.py}
  \item Sweep xea: \texttt{examples/ABRAMSON1965/run\_uncoupled\_xea\_xcm\_sweep.py}
  \item Baseline: \texttt{examples/ABRAMSON1965/run\_baseline\_sensitivity.py}
  \item Confronto $\mathrm{Im}(Q)$: \texttt{COUPLING/hydroelastic\_utl/abramson\_imq\_comparison.py}
  \item Output root: \texttt{output\_data/ABRAMSON1965/}
  \item Capytaine per caso: \texttt{FLUID/capytaine/ABRAMSON1965/}
\end{itemize}

\end{document}
